\documentclass[conference]{IEEEtran}
\usepackage[utf8]{inputenc}
\usepackage[ukrainian]{babel}
\usepackage{graphicx} % Пакет для роботи з графікою
\usepackage{amsmath}
\usepackage{caption}

\title{Наукова візуалізація та графічна інтерпретація результатів дослідження успішності студентів}
\author{Дмитро Страхов, Спеціальність 126}

\begin{document}

\maketitle

\begin{abstract}
У даній роботі представлено результати візуального аналізу навчальних показників. Використано методи статистичної візуалізації (Distribution plots, Box-plots, Heatmaps) для виявлення закономірностей у даних.
\end{abstract}

\section{Вступ}
Для прийняття рішень у сучасних ІС важливо не лише розраховувати статистики, а й правильно їх візуалізувати. Ми використовуємо бібліотеки Seaborn та Plotly для побудови графіків журнальної якості.

\section{Аналіз розподілу даних}
На рис. \ref{fig:dist} представлено гістограму розподілу підсумкових балів. Використання кривої щільності (KDE) дозволяє оцінити форму розподілу та наявність аномалій.

\begin{figure}[h]
    \centering
    \includegraphics[width=0.9\linewidth]{distribution_plot.pdf}
    \caption{Гістограма розподілу total\_score з накладеною кривою KDE.}
    \label{fig:dist}
\end{figure}

\section{Порівняльний аналіз та кореляція}
\subsection{Діаграми "ящик з вусами"}
Для порівняння груп студентів за рівнем відвідуваності побудовано Box-plot (див. рис. \ref{fig:box}). Це дозволяє ідентифікувати медіану та міжквартильний розмах.

\begin{figure}[h]
    \centering
    \includegraphics[width=0.85\linewidth]{boxplot_attendance.png}
    \caption{Порівняння успішності студентів з різним рівнем відвідуваності.}
    \label{fig:box}
\end{figure}

\subsection{Кореляційна матриця}
На рис. \ref{fig:heat} наведено теплову карту, яка візуалізує силу зв'язку між параметрами. Коефіцієнт кореляції Пірсона відображений у кожній комірці.

\begin{figure}[h]
    \centering
    \includegraphics[width=1\linewidth]{correlation_heatmap.pdf}
    \caption{Теплова карта кореляції навчальних метрик.}
    \label{fig:heat}
\end{figure}

\section{Висновки}
Використання векторних форматів (PDF/SVG) забезпечує високу чіткість графіків при масштабуванні, що є стандартом для наукових видань Q1.

\section{Методи (Methods) — Побудова моделі}
Для прогнозування успішності використано модель множинної лінійної регресії:
\begin{equation}
    y = \beta_0 + \beta_1 x_1 + \beta_2 x_2 + \beta_3 x_3 + \epsilon
\end{equation}
де $y$ — прогнозований бал (total\_score), $x_1$ — години навчання, $x_2$ — відвідуваність.

Для оцінки якості моделі використано коефіцієнт детермінації $R^2$:
\begin{equation}
    R^2 = 1 - \frac{\sum (y_i - \hat{y}_i)^2}{\sum (y_i - \bar{y}_i)^2}
\end{equation}
Значення $R^2$ показує частку дисперсії залежної змінної, що пояснюється моделлю.

\section{Результати та оцінка точності (Results and Metrics)}
Для перевірки точності прогнозування було розраховано основні метрики регресії. 
Середньоквадратична помилка (MSE) склала [впиши своє число], що вказує на відносно низьку дисперсію залишків. 

Коефіцієнт детермінації $R^2 = [впиши своє число]$ підтверджує, що вхідні параметри (години навчання та відвідуваність) мають визначальний вплив на цільову змінну. 

\section{Результати та оцінка точності (Results and Metrics)}
Для перевірки точності прогнозування було розраховано основні метрики регресії. 
Середньоквадратична помилка (MSE) склала [впиши своє число], що вказує на відносно низьку дисперсію залишків. 

Коефіцієнт детермінації $R^2 = [впиши своє число]$ підтверджує, що вхідні параметри (години навчання та відвідуваність) мають визначальний вплив на цільову змінну. 
На рис. \ref{fig:reg} наведено графік залежності прогнозованих значень від реальних, де червона лінія позначає ідеальну відповідність.
\end{document}
\begin{document}

\maketitle

\section{Introduction}

\end{document}
